All times are the mean time of one pass over 10{,}000 images, in milliseconds (ms), from \val{Sp:runs} recorded passes per configuration on the Spartan node unless the Apple M4 is named. Batch sizes are in images per batch; ``$S$'' is the speedup over the serial code on the same machine and batch size (in multiples, $\times$), and ``$E$'' is the efficiency $S/p$. Compute-only comparisons (stages 1--3, no preprocessing and no copies) and end-to-end comparisons (stages 0--3 including copies) are reported separately. The M4 is used only for the OpenMP study in Section~\ref{sec:omp}; no CPU-versus-GPU comparison uses it.

\subsection{OpenMP Scalability (RQ1)}\label{sec:omp}

On the Xeon the speedup is close to linear for batches of 32 images or more (Figure~\ref{fig:omp}, Table~\ref{tab:omp}). With 24 threads the speedup is \val{Sp:omp:speedup:4096:24}$\times$ at batch 4096 (efficiency \val{Sp:omp:eff:4096:24}) and \val{Sp:omp:speedup:256:24}$\times$ at batch 256 (efficiency \val{Sp:omp:eff:256:24}). With one image per batch the best result is \val{Sp:omp:best:speedup:1}$\times$ with \val{Sp:omp:best:thr:1} threads, and the efficiency at 24 threads falls to \val{Sp:omp:eff:1:24}: a batch of one image offers only 16 work items per stage and each stage ends in a barrier. OpenMP with one thread needs \val{Sp:omp:e2e:256:1}\,ms against \val{Sp:serial:e2e:256}\,ms for the serial code (batch 256), so the OpenMP runtime adds no measurable cost.

The Apple M4 behaves differently. The speedup follows the ideal line up to its four performance cores (\val{Mac:omp:speedup:4096:4}$\times$ with 4 threads at batch 4096) and then flattens to \val{Mac:omp:speedup:4096:10}$\times$ with 10 threads (efficiency \val{Mac:omp:eff:4096:10}); the dotted line in Figure~\ref{fig:omp} marks the four performance cores. With one image per batch OpenMP is \emph{slower} than serial on the M4 for every thread count above one: \val{Mac:omp:speedup:1:2}$\times$ with 2 threads and \val{Mac:omp:speedup:1:10}$\times$ with 10 threads. The serial pass takes \val{Mac:serial:e2e:4096}\,ms on the M4 and \val{Sp:serial:e2e:4096}\,ms on one Xeon core (a CPU-to-CPU comparison; the build settings differ, Table~\ref{tab:compiler}).

\widefigure{omp_combined.pdf}{OpenMP speedup over serial (top) and parallel efficiency (bottom) versus thread count, for four batch sizes, on the Apple M4 (left) and the Xeon (right). The dashed lines are ideal scaling; the dotted line marks the four performance cores of the M4.}{fig:omp}

\begin{widetable}{OpenMP scaling on the Spartan node: time per pass (ms), speedup $S$ over serial ($\times$) and efficiency $E=S/p$ for three batch sizes.}{tab:omp}
\begin{tabular}{rrrrrrrrrr}
\toprule
Threads & \multicolumn{3}{c}{Batch 1} & \multicolumn{3}{c}{Batch 256} & \multicolumn{3}{c}{Batch 4096} \\
\cmidrule(lr){2-4}\cmidrule(lr){5-7}\cmidrule(lr){8-10}
 & Time (ms) & $S$ ($\times$) & $E$ & Time (ms) & $S$ ($\times$) & $E$ & Time (ms) & $S$ ($\times$) & $E$ \\
\midrule
1 & 1{,}221.8 & 1.01 & 1.01 & 1{,}205.5 & 1.02 & 1.02 & 1{,}205.5 & 1.02 & 1.02 \\
2 & 662.6 & 1.86 & 0.93 & 631.0 & 1.96 & 0.98 & 629.5 & 1.96 & 0.98 \\
4 & 370.8 & 3.33 & 0.83 & 333.8 & 3.70 & 0.93 & 334.0 & 3.70 & 0.92 \\
8 & 212.6 & 5.81 & 0.73 & 168.6 & 7.32 & 0.92 & 166.8 & 7.41 & 0.93 \\
12 & 216.1 & 5.72 & 0.48 & 114.7 & 10.77 & 0.90 & 112.6 & 10.98 & 0.91 \\
16 & 162.8 & 7.59 & 0.47 & 90.3 & 13.68 & 0.86 & 85.5 & 14.45 & 0.90 \\
24 & 167.8 & 7.36 & 0.31 & 58.9 & 20.99 & 0.87 & 57.3 & 21.55 & 0.90 \\
\bottomrule
\end{tabular}

\end{widetable}

\subsection{End-to-End Comparison}\label{sec:e2e}

This is the \emph{end-to-end} comparison of the four primary implementations (stages 0--3, copies included; the wall-clock time of a pass). Table~\ref{tab:cpugpu} and Figures~\ref{fig:exec}--\ref{fig:tput} show the results. The CUDA entries use a fixed 16$\times$16 block and one preprocessing thread, whereas OpenMP and the hybrid use their best thread count, so the comparison is ``CUDA with a fixed block against OpenMP with tuned thread count''. The serial and OpenMP times hardly depend on the batch size: the serial pass takes \val{Sp:serial:e2e:1}\,ms with batch 1 and \val{Sp:serial:e2e:4096}\,ms with batch 4096, because every image is processed independently by the same loops. The GPU depends strongly on the batch size: \cudabase{} falls from \val{Sp:cuda:e2e:1}\,ms (batch 1) to \val{Sp:cuda:e2e:4096}\,ms (batch 4096).

With one image per batch the GPU pipeline is slower than serial (\val{Sp:cuda:speedup:1}$\times$): each of the 10{,}000 batches needs three kernel launches, two copies and a synchronization, and the kernels alone take \val{Sp:cuda:kernel:1}\,ms per pass, \val{Sp:cuda:kus:1}\,$\mu$s per batch. From batch 8 on it is faster than serial, reaching \val{Sp:cuda:speedup:4096}$\times$ end-to-end at batch 4096.

Against the best OpenMP configuration, the CUDA pipeline is slower for batches up to 64, level at batch 128 (\val{Sp:cuda:e2e:128}\,ms against \val{Sp:omp:best:e2e:128}\,ms) and ahead from batch 256 (\val{Sp:cuda:e2e:256}\,ms against \val{Sp:omp:best:e2e:256}\,ms); at batch 4096 it is only \val{Sp:cuda:vscpu:4096}$\times$ faster than 24-thread OpenMP. Using the rate convention of Section~\ref{sec:method} (stages 1--2 only, divided by the end-to-end time unless marked kernel-only), a pass at batch 4096 reaches \val{Sp:serial:gflops:4096}\,GFLOP/s serial, \val{Sp:omp:best:gflops:4096}\,GFLOP/s with OpenMP and \val{Sp:cuda:gflops:4096}\,GFLOP/s for the CUDA pipeline end-to-end; the kernel-only rate of the CUDA kernels is \val{Sp:cuda:kgflops:4096}\,GFLOP/s.

Against the same baselines the pipelined hybrid is the fastest implementation from batch \val{Sp:hyb:crossover} on, \val{Sp:hyb:vscpu:4096}$\times$ faster than the best OpenMP configuration at batch 4096, whereas OpenMP is fastest for batches of 1 and 8 (Section~\ref{sec:hybrid}). Figure~\ref{fig:speedup} shows the speedup over serial of the four implementations and Figure~\ref{fig:tput} the throughput in images per second. At batch 4096 the serial code processes \val{Sp:serial:tput:4096}~images/s, OpenMP \val{Sp:omp:best:tput:4096}~images/s, the CUDA implementation with single-threaded preprocessing \val{Sp:cuda:tput:4096}~images/s and the pipelined hybrid \val{Sp:hyb:best:tput:4096}~images/s. Speedup and throughput are two views of the same time measurements: $S=T_{\mathrm{serial}}/T$ and throughput $=10{,}000/T$.

\begin{widetable}{End-to-end comparison on the Spartan node: mean time per pass (ms) and speedup over serial ($\times$). OpenMP and Hybrid show the best of the tested thread counts (in brackets). ``CUDA end-to-end'' preprocesses on one CPU thread; the kernel-only time of Table~\ref{tab:compute} excludes preprocessing and copies. $^{\ast}$The kernel-only speedup is a compute-only ratio (serial time of stages 1--3 divided by the kernel time); all other speedups are end-to-end ratios.}{tab:cpugpu}
\begin{tabular}{rrrrrrrrr}
\toprule
Batch & Serial & OpenMP (threads) & CUDA kernel & CUDA end-to-end & Hybrid (threads) & \multicolumn{3}{c}{Speedup over serial ($\times$)} \\
\cmidrule(lr){7-9}
 & (ms) & (ms) & only (ms) & (ms) & (ms) & kernel-only & CUDA e2e & Hybrid \\
\midrule
1 & 1{,}234.8 & 162.8 (16) & 1{,}091.0 & 1{,}507.0 & 650.9 (1) & 1.13 & 0.82 & 1.90 \\
8 & 1{,}235.7 & 75.5 (24) & 143.9 & 218.8 & 87.5 (8) & 8.59 & 5.65 & 14.13 \\
32 & 1{,}235.0 & 61.9 (24) & 47.3 & 85.5 & 33.3 (2) & 26.10 & 14.45 & 37.09 \\
64 & 1{,}234.8 & 59.2 (24) & 34.4 & 65.9 & 27.9 (24) & 35.87 & 18.73 & 44.20 \\
128 & 1{,}234.6 & 58.2 (24) & 29.8 & 58.6 & 26.5 (24) & 41.43 & 21.06 & 46.59 \\
256 & 1{,}235.2 & 58.9 (24) & 28.1 & 55.1 & 26.0 (24) & 44.01 & 22.40 & 47.54 \\
1024 & 1{,}235.0 & 57.7 (24) & 23.7 & 48.3 & 25.2 (24) & 52.02 & 25.56 & 48.98 \\
4096 & 1{,}235.4 & 57.3 (24) & 23.0 & 46.7 & 24.7 (24) & 53.65 & 26.44 & 50.02 \\
\bottomrule
\end{tabular}

\end{widetable}

\colfigure{exec_time_spartan_k40_native.pdf}{End-to-end time per pass (ms, log scale) versus batch size for the four primary implementations (end-to-end). Error bars are 95\,\% confidence intervals (smaller than the markers).}{fig:exec}

\colfigure{speedup_spartan_k40_native.pdf}{End-to-end speedup over the serial CPU ($T_{\mathrm{serial}}/T$, same batch size) of OpenMP (best thread count), the CUDA implementation with single-threaded preprocessing and the pipelined hybrid. The serial CPU is the baseline at 1.0.}{fig:speedup}

\colfigure{throughput_spartan_k40_native.pdf}{Throughput (images per second, log scale) of the four implementations versus batch size, end-to-end ($10{,}000/T$).}{fig:tput}

\subsection{CUDA Kernel Versus End-to-End Performance}\label{sec:compute}

This is the \emph{compute-only} comparison. Table~\ref{tab:compute} compares only the computation of the hidden layer, the output layer and the argmax; preprocessing and copies are excluded, so the GPU is represented by its kernels alone. For the CPU it is the sum of the timed stage durations. The GPU kernels are \val{Sp:comp:gpuvsserial:4096}$\times$ faster than serial at batch 4096 (\val{Sp:comp:kernel:4096}\,ms against \val{Sp:comp:serial:4096}\,ms) and \val{Sp:comp:gpuvsomp:4096}$\times$ faster than the best OpenMP configuration (\val{Sp:comp:omp:4096}\,ms). They beat OpenMP from batch 32 on (\val{Sp:comp:gpuvsomp:32}$\times$ at batch 32). For one image per batch the kernels are only \val{Sp:comp:gpuvsserial:1}$\times$ faster than serial and \val{Sp:comp:gpuvsomp:1}$\times$ as fast as OpenMP, i.e.\ slower: the batch gives a GPU too little parallel work.

Figure~\ref{fig:kernel} contrasts the two views for the CUDA implementation. The kernel-only speedup over the serial stages 1--3 is \val{Sp:comp:gpuvsserial:4096}$\times$ at batch 4096, about twice the \emph{end-to-end} speedup of \val{Sp:cuda:speedup:4096}$\times$, because the end-to-end time also contains the single-threaded preprocessing, the copies and the launch overhead (Section~\ref{sec:overhead}). The two ratios answer different questions and must not be mixed: the first compares only the neural-network computation, the second the complete pipeline an application waits for.

\begin{widetable}{Compute-only comparison on the Spartan node (stages 1--3, no preprocessing, no copies): time per pass (ms). OpenMP shows the best thread count in brackets; ``GPU kernels'' is the CUDA-event time of the three kernels (16$\times$16 block).}{tab:compute}
\begin{tabular}{rrrrrr}
\toprule
Batch & Serial & OpenMP (threads) & GPU kernels & \multicolumn{2}{c}{GPU kernels' speedup ($\times$)} \\
\cmidrule(lr){5-6}
 & (ms) & (ms) & (ms) & vs serial & vs OpenMP \\
\midrule
1 & 1{,}213.7 & 126.4 (16) & 1{,}091.0 & 1.11 & 0.12 \\
8 & 1{,}215.4 & 70.5 (24) & 143.9 & 8.45 & 0.49 \\
32 & 1{,}214.8 & 60.0 (24) & 47.3 & 25.67 & 1.27 \\
64 & 1{,}214.6 & 57.9 (24) & 34.4 & 35.28 & 1.68 \\
128 & 1{,}214.4 & 57.0 (24) & 29.8 & 40.75 & 1.91 \\
256 & 1{,}214.9 & 57.8 (24) & 28.1 & 43.29 & 2.06 \\
1024 & 1{,}214.8 & 56.7 (24) & 23.7 & 51.17 & 2.39 \\
4096 & 1{,}215.2 & 56.1 (24) & 23.0 & 52.77 & 2.44 \\
\bottomrule
\end{tabular}

\end{widetable}

\widefigure{cuda_kernel_vs_e2e_spartan_k40_native.pdf}{\cudabase{}: kernel-only time against end-to-end time (left, ms) and the corresponding speedups (right, $\times$): kernel-only against the serial time of stages 1--3 (compute-only), end-to-end against the serial end-to-end time. Kernel-only numbers exclude preprocessing, copies and launch overhead.}{fig:kernel}

\subsection{CUDA Transfer and Preprocessing Overhead (RQ4)}\label{sec:overhead}

Table~\ref{tab:breakdown} and Figure~\ref{fig:breakdown} split the pass of the CUDA implementation into stages. The two copies take \val{Sp:cuda:tshare:1}\,\% of the end-to-end time at batch 1 and \val{Sp:cuda:tshare:4096}\,\% at batch 4096. The host-to-device copy of the 31.4\,MB takes \val{Sp:cuda:h2d:4096}\,ms at batch 4096, which is \val{Sp:cuda:h2dgbs:4096}\,GB/s and agrees with the \val{Sp:micro:h2d_pageable_bandwidth:0}\,GB/s measured separately for pageable memory (Table~\ref{tab:micro}); pinned memory reaches \val{Sp:micro:h2d_pinned_bandwidth:0}\,GB/s.

At large batches the bigger non-GPU cost is the single-threaded preprocessing, which takes a constant \val{Sp:cuda:pre:4096}\,ms per pass. Its share of the end-to-end time grows from \val{Sp:cuda:preshare:1}\,\% at batch 1 to \val{Sp:cuda:preshare:4096}\,\% at batch 4096, while the kernels fall from \val{Sp:cuda:kshare:1}\,\% to \val{Sp:cuda:kshare:4096}\,\%. At batch 1 the launch, API and synchronization overhead (column ``Other'') is \val{Sp:cuda:otherms:1}\,ms per pass, \val{Sp:cuda:othershare:1}\,\% of the total, and it vanishes for large batches. This is why the end-to-end speedup stays well below the kernel-only speedup.

\colfigure{transfer_breakdown_spartan_k40_native.pdf}{Share of the end-to-end time of the CUDA implementation (single-threaded preprocessing) spent in each stage (\%) versus batch size.}{fig:breakdown}

\begin{widetable}{Time per pass of the CUDA implementation by stage on the Spartan node (ms). ``Other'' is the part of the wall-clock time not covered by the four timed stages (launch, API calls, synchronization).}{tab:breakdown}
\begin{tabular}{rrrrrrrr}
\toprule
Batch & Preprocess & H2D copy & Kernels & D2H copy & Other & Total & Copies share \\
 & (ms) & (ms) & (ms) & (ms) & (ms) & (ms) & (\% of total) \\
\midrule
1 & 20.4 & 85.4 & 1{,}091.0 & 213.3 & 96.8 & 1{,}507.0 & 19.8 \\
8 & 20.2 & 17.2 & 143.9 & 25.9 & 11.6 & 218.8 & 19.7 \\
32 & 20.2 & 7.8 & 47.3 & 7.0 & 3.1 & 85.5 & 17.3 \\
64 & 20.2 & 6.6 & 34.4 & 3.1 & 1.6 & 65.9 & 14.7 \\
128 & 20.2 & 6.0 & 29.8 & 1.8 & 0.8 & 58.6 & 13.2 \\
256 & 20.2 & 5.5 & 28.1 & 0.9 & 0.4 & 55.1 & 11.7 \\
1024 & 20.3 & 3.8 & 23.7 & 0.3 & 0.1 & 48.3 & 8.6 \\
4096 & 20.2 & 3.3 & 23.0 & 0.1 & 0.0 & 46.7 & 7.3 \\
\bottomrule
\end{tabular}

\end{widetable}

\begin{coltable}{Platform limits measured by micro-benchmarks.}{tab:micro}
\begin{tabular}{@{}lr@{}}
\toprule
Quantity & Measured value \\
\midrule
PCIe host-to-device, pageable / pinned & 9.5 / 10.5 GB/s \\
PCIe device-to-host, pageable / pinned & 10.0 / 10.5 GB/s \\
GPU device-to-device copy & 189.5 GB/s \\
Kernel launch, asynchronous / with sync & 3.12 / 7.44 $\mu$s \\
Xeon memory bandwidth, 1 / 8 / 24 threads & 11.8 / 46.5 / 43.4 GB/s \\
Xeon OpenMP fork--join, 8 / 24 threads & 2.13 / 7.16 $\mu$s \\
M4 memory bandwidth, 1 / 10 threads & 90.0 / 86.5 GB/s \\
M4 OpenMP fork--join, 4 / 10 threads & 7.25 / 21.98 $\mu$s \\
\bottomrule
\end{tabular}

\end{coltable}

\subsection{Hybrid OpenMP--CUDA Performance (RQ5)}\label{sec:hybrid}

Figure~\ref{fig:hybrid} and Table~\ref{tab:hybrid} compare OpenMP, the CUDA implementation, the hybrid-sequential control and the pipelined hybrid. The pipelined hybrid is faster than \cudabase{} at every batch size: \val{Sp:hyb:vscuda:1}$\times$ at batch 1, \val{Sp:hyb:vscuda:32}$\times$ at batch 32 and \val{Sp:hyb:vscuda:4096}$\times$ at batch 4096. Welch's $t$-test gives $p\le$ \val{Sp:test:hyb-cuda:maxp} for every batch size, and the 95\,\% bootstrap interval of the time ratio at batch 4096 is \val{Sp:test:hyb-cuda:cilo:4096} to \val{Sp:test:hyb-cuda:cihi:4096}. Against the best OpenMP configuration the hybrid is faster from batch \val{Sp:hyb:crossover} on, \val{Sp:hyb:vscpu:4096}$\times$ at batch 4096. At batch 1 and batch 8 OpenMP alone is faster (\val{Sp:omp:best:e2e:1}\,ms against \val{Sp:hyb:best:e2e:1}\,ms at batch 1). The CUDA pipeline with single-threaded preprocessing was not the fastest implementation at any batch size (Table~\ref{tab:best}).

Two ingredients contribute to the hybrid gain. Parallel preprocessing alone (the hybrid-sequential control) gives \val{Sp:hseq:vscuda:4096}$\times$ over \cudabase{} at batch 4096, and overlapping the batches adds another \val{Sp:hyb:vshseq:4096}$\times$. Even with a single preprocessing thread, pipelining with pinned memory gives \val{Sp:hyb:vscuda1:4096}$\times$. At batch 4096 the hybrid pass takes \val{Sp:hyb:best:e2e:4096}\,ms, only \val{Sp:hyb:vskernel:4096}$\times$ the kernel-only time of the CUDA implementation (\val{Sp:cuda:kernel:4096}\,ms): preprocessing and copies are almost completely hidden behind the kernels of other batches, so the GPU computation is again the bottleneck. Because the GPU stages overlap in the pipelined hybrid, only its end-to-end time is meaningful.

\widefigure{hybrid_compare_spartan_k40_native.pdf}{Left: end-to-end time per pass (ms, log scale) of OpenMP, the CUDA implementation with single-threaded preprocessing, the hybrid-sequential control and the pipelined hybrid (best thread counts). Right: speedup of the control and of the pipelined hybrid over the CUDA implementation ($\times$).}{fig:hybrid}

\begin{widetable}{Hybrid ablation on the Spartan node: time per pass (ms), speedup ($\times$) and significance. The first column is the CUDA implementation with one preprocessing thread; the hybrid-sequential column is a diagnostic control. Thread counts of the best configuration are in brackets.}{tab:hybrid}
\fitwide{\begin{tabular}{rrrrrrrr}
\toprule
Batch & CUDA-only & Hybrid-sequential (thr.) & Hybrid-pipelined (thr.) & Seq. vs CUDA-only & Pipelined vs CUDA-only & Pipelined vs OpenMP & Welch $p$ \\
 & (ms) & (ms) & (ms) & ($\times$) & ($\times$) & ($\times$) & (pipelined vs CUDA-only) \\
\midrule
1 & 1{,}507.0 & 1{,}520.6 (1) & 650.9 (1) & 0.99 & 2.32 & 0.25 & $8.7\times10^{-61}$ \\
8 & 218.8 & 207.5 (8) & 87.5 (8) & 1.05 & 2.50 & 0.86 & $2.4\times10^{-103}$ \\
32 & 85.5 & 69.9 (12) & 33.3 (2) & 1.22 & 2.57 & 1.86 & $3.9\times10^{-55}$ \\
64 & 65.9 & 49.9 (16) & 27.9 (24) & 1.32 & 2.36 & 2.12 & $1.6\times10^{-69}$ \\
128 & 58.6 & 41.9 (16) & 26.5 (24) & 1.40 & 2.21 & 2.20 & $1.2\times10^{-77}$ \\
256 & 55.1 & 38.2 (24) & 26.0 (24) & 1.44 & 2.12 & 2.27 & $5.9\times10^{-79}$ \\
1024 & 48.3 & 31.3 (12) & 25.2 (24) & 1.55 & 1.92 & 2.29 & $4.6\times10^{-33}$ \\
4096 & 46.7 & 28.8 (16) & 24.7 (24) & 1.62 & 1.89 & 2.32 & $4.8\times10^{-62}$ \\
\bottomrule
\end{tabular}
}
\end{widetable}

\begin{coltable}{Fastest implementation per batch size (end-to-end time per pass in ms; gap = second fastest / fastest). ``CUDA (1-thr. pre.)'' is the CUDA implementation with single-threaded preprocessing and a fixed 16$\times$16 block.}{tab:best}
\fitcol{\begin{tabular}{rlrlrr}
\toprule
Batch & Fastest & Time & Second fastest & Time & Gap \\
 & implementation & (ms) & implementation & (ms) & ($\times$) \\
\midrule
1 & OpenMP & 162.8 & Hybrid-pipelined & 650.9 & 4.00 \\
8 & OpenMP & 75.5 & Hybrid-pipelined & 87.5 & 1.16 \\
32 & Hybrid-pipelined & 33.3 & OpenMP & 61.9 & 1.86 \\
64 & Hybrid-pipelined & 27.9 & OpenMP & 59.2 & 2.12 \\
128 & Hybrid-pipelined & 26.5 & OpenMP & 58.2 & 2.20 \\
256 & Hybrid-pipelined & 26.0 & CUDA (1-thr. pre.) & 55.1 & 2.12 \\
1024 & Hybrid-pipelined & 25.2 & CUDA (1-thr. pre.) & 48.3 & 1.92 \\
4096 & Hybrid-pipelined & 24.7 & CUDA (1-thr. pre.) & 46.7 & 1.89 \\
\bottomrule
\end{tabular}
}
\end{coltable}

\subsection{CUDA Block Size and Tuning}\label{sec:block}

The block size of the hidden-layer kernel matters only for small batches (Table~\ref{tab:block}, Figure~\ref{fig:block}). At batch 64 the kernels take \val{Sp:cuda:blk8:kernel:64}\,ms with 8$\times$8 blocks, \val{Sp:cuda:blk16:kernel:64}\,ms with 16$\times$16 and \val{Sp:cuda:blk32:kernel:64}\,ms with 32$\times$32 blocks. At batch 4096 the three differ by less than 1\,\% (\val{Sp:cuda:blk8:kernel:4096}, \val{Sp:cuda:blk16:kernel:4096} and \val{Sp:cuda:blk32:kernel:4096}\,ms). A 32$\times$32 block at batch 32 launches only 8 thread blocks for the 15 multiprocessors of the K40m, consistent with its slower result (\val{Sp:cuda:blk32:kernel:32}\,ms against \val{Sp:cuda:blk8:kernel:32}\,ms for 8$\times$8). If the best of the three block sizes were used for the end-to-end time of the CUDA implementation, it would improve by at most \val{Sp:cuda:blockgain:max}\,\% (at batch \val{Sp:cuda:blockgain:argb}), and the batch size at which it catches up with the best OpenMP configuration would move from 256 to \val{Sp:cuda:bestblock:crossover} (where the two are within 0.1\,\%). The main comparison therefore keeps the fixed 16$\times$16 block and is labelled accordingly; it is not a best-versus-best comparison, because the hybrid was not block-tuned either.

\begin{widetable}{CUDA block-size comparison on the Spartan node: time per pass (ms) for the hidden-layer block sizes 8$\times$8, 16$\times$16 and 32$\times$32 (kernels only and end-to-end), the fastest block for the end-to-end time, and its gain over the fixed 16$\times$16 block. Block sizes of the output-layer and argmax kernels are fixed at 128 threads.}{tab:block}
\begin{tabular}{rrrrrrrcr}
\toprule
Batch & \multicolumn{3}{c}{CUDA kernels only (ms)} & \multicolumn{3}{c}{CUDA end-to-end (ms)} & Fastest & Gain over \\
\cmidrule(lr){2-4}\cmidrule(lr){5-7}
 & 8$\times$8 & 16$\times$16 & 32$\times$32 & 8$\times$8 & 16$\times$16 & 32$\times$32 & block (end-to-end) & 16$\times$16 (\%) \\
\midrule
1 & 1{,}095.6 & 1{,}091.0 & 1{,}136.9 & 1{,}510.1 & 1{,}507.0 & 1{,}560.6 & 16$\times$16 & 0.0 \\
8 & 142.7 & 143.9 & 144.8 & 218.3 & 218.8 & 219.6 & 8$\times$8 & 0.3 \\
32 & 45.3 & 47.3 & 56.7 & 83.2 & 85.5 & 94.8 & 8$\times$8 & 2.7 \\
64 & 37.3 & 34.4 & 48.7 & 69.1 & 65.9 & 81.1 & 16$\times$16 & 0.0 \\
128 & 29.5 & 29.8 & 34.6 & 58.2 & 58.6 & 63.3 & 8$\times$8 & 0.7 \\
256 & 26.5 & 28.1 & 29.6 & 53.6 & 55.1 & 56.6 & 8$\times$8 & 2.8 \\
1024 & 23.6 & 23.7 & 24.2 & 48.3 & 48.3 & 48.8 & 8$\times$8 & 0.0 \\
4096 & 23.1 & 23.0 & 23.2 & 46.8 & 46.7 & 46.8 & 16$\times$16 & 0.0 \\
\bottomrule
\end{tabular}

\end{widetable}

\colfigure{cuda_block_size_spartan_k40_native.pdf}{CUDA kernel time per pass (ms, log scale) for the thread-block sizes 8$\times$8, 16$\times$16 and 32$\times$32 of the hidden-layer kernel.}{fig:block}

\subsection{Initialization and Cold Start}\label{sec:cold}

The one-time costs that are excluded from the steady-state times were measured separately in 10 fresh processes (batch 256, 24 threads; Table~\ref{tab:cold} and Figure~\ref{fig:cold}). Creating the CUDA context takes \val{Sp:cold:init:mean}\,ms on average (median \val{Sp:cold:init:median}\,ms, range \val{Sp:cold:init:min}--\val{Sp:cold:init:max}\,ms), the weight upload \val{Sp:cold:upload:mean}\,ms and the buffer, stream and event allocation \val{Sp:cold:alloc:mean}\,ms, in total \val{Sp:cold:ctor:mean}\,ms, which is \val{Sp:cold:ctorvspass}$\times$ the time of one steady-state pass of the CUDA implementation. After this set-up the first CUDA pass takes \val{Sp:cold:cuda1:mean}\,ms and later passes \val{Sp:cold:cuda3:mean}\,ms, so the GPU kernels show no first-pass penalty. The first pipelined-hybrid pass is \val{Sp:cold:hyb:firstpenalty}\,\% slower than the second, and the first 24-thread OpenMP pass is \val{Sp:cold:omp:firstpenalty}\,\% slower (\val{Sp:cold:omp1:mean}\,ms against \val{Sp:cold:omp2:mean}\,ms); the serial code shows no difference. A program that processes a single pass of 10{,}000 images would therefore pay about \val{Sp:cold:ctor:mean}\,ms of GPU initialization on top of its pass; the steady-state numbers apply to programs that amortize this cost.

\widefigure{cold_start.pdf}{Cold start versus steady state on the Spartan node (10 fresh processes, batch 256, 24 OpenMP threads, mean with standard deviation). (a)~Time of the first and the second pass after the one-time GPU set-up. (b)~Components of the one-time GPU set-up, which is excluded from all steady-state times, against one steady-state CUDA pass.}{fig:cold}

\begin{widetable}{Cold-start measurements on the Spartan node (10 fresh processes, batch 256, 24 OpenMP threads): one-time GPU set-up and the first passes against later passes.}{tab:cold}
\begin{tabular}{lrrrrr}
\toprule
Quantity & Mean & Median & Std. dev. & Min & Max \\
 & (ms) & (ms) & (ms) & (ms) & (ms) \\
\midrule
CUDA context creation & 106.83 & 104.64 & 8.11 & 103.19 & 129.79 \\
Weight upload & 0.40 & 0.40 & 0.01 & 0.38 & 0.42 \\
Buffer, stream and event allocation & 2.97 & 2.94 & 0.08 & 2.88 & 3.11 \\
Total one-time GPU set-up & 110.65 & 108.53 & 8.07 & 106.93 & 133.49 \\
CUDA pass, 1st & 54.99 & 54.96 & 0.24 & 54.67 & 55.35 \\
CUDA pass, 3rd & 55.02 & 55.05 & 0.12 & 54.77 & 55.17 \\
Hybrid pass, 1st & 27.35 & 27.47 & 0.36 & 26.79 & 27.78 \\
Hybrid pass, 2nd & 25.95 & 25.93 & 0.12 & 25.80 & 26.19 \\
OpenMP (24 threads) pass, 1st & 80.84 & 80.45 & 1.55 & 79.90 & 85.20 \\
OpenMP (24 threads) pass, 2nd & 63.22 & 63.25 & 0.23 & 62.75 & 63.49 \\
Serial pass, 1st & 1{,}231.70 & 1{,}231.46 & 1.13 & 1{,}230.34 & 1{,}234.09 \\
Serial pass, 2nd & 1{,}232.32 & 1{,}232.13 & 0.96 & 1{,}231.11 & 1{,}234.06 \\
\bottomrule
\end{tabular}

\end{widetable}

